% =====================================================================
% 第 9 节：固定绝对初值下边界的“精确/近似”表述修订（逐段替换片段）
% 基于 GitHub 当前版本 436dedd 的 latex/flatten_curve_analysis_cn.tex
% 行号均指原稿行号。同样内容的 git 补丁见 sec9_boundary.patch。
% 不新增编号公式、定理或标签，现有公式与定理编号不变。
% =====================================================================


% ---------------------------------------------------------------------
% 【改动 1】§9.1 设定与比较基准：替换原第 1515–1517 行
%   原文开头：“数值计算固定绝对初值 $I_0=1.00659\times10^{-3}$，不随……”
%   至：“……不给出所有初值或人口下统一的误差认证。”
% ---------------------------------------------------------------------
数值计算在各人口情景下沿用全市人口下拟合得到的绝对初值 $I_0=1.00659\times10^{-3}$，并取 $S_0=N-I_0$，故 $i_0=I_0/N$ 与 $s_0=1-i_0$ 随人口规模改变。这是本节采用的情景比较约定；若在不同人口下重新拟合，所得初值未必相同。第\ref{sec:scaling-theory}节的规模不变性以固定归一化初值为前提，两者需加以区分。

固定绝对初值时，成本条件对应的阈值比例应视为 $i_0$ 的函数 $\theta_{\rm cost}(i_0)$，由 $J(\theta,i_0)=J^{\rm T}$ 确定，其中 $s_0=1-i_0$ 随 $i_0$ 同步改变。$(N,\eta)$ 平面上的精确成本边界为
\[
\eta_{\rm cost}(N)=N\,\theta_{\rm cost}(I_0/N),
\]
时长边界同理。这里的 $\theta_{\rm cost}(0)$ 指 $i_0\downarrow0$ 时由正初值的首次触发方程和控制期指标公式确定的极限，不对应 $I_0=0$ 的疫情：后者有 $I\equiv0$，不会触发控制。在内部触发区间 $s^*>s_c$ 上，首次触发方程关于 $s^*$ 的导数非零，$s^*(\theta,i_0)$ 可光滑延拓至 $i_0=0$，代入控制期公式后 $J,\Delta t$ 关于 $(\theta,i_0)$ 二阶连续可微；命题\ref{prop:threshold-tradeoff}证明中两项的符号在延拓后不变，故 $\partial J/\partial\theta<0$。由隐函数定理，
\[
\eta_{\rm cost}(N)=N\,\theta_{\rm cost}(0)+\kappa_{\rm cost}I_0+O\!\left(\frac{I_0^2}{N}\right),
\qquad
\kappa_{\rm cost}=-\left.\frac{\partial J/\partial i_0}{\partial J/\partial\theta}\right|_{(\theta_{\rm cost}(0),\,0)}.
\]
对首次触发方程隐式求导可得 $\partial s^*/\partial i_0>0$，故固定 $\theta$ 时 $J$ 随 $i_0$ 增加，$\kappa_{\rm cost}>0$。因此，精确边界一般不是直线；当 $I_0/N$ 很小时，其一阶近似是斜率为 $\theta_{\rm cost}(0)$、截距为 $\kappa_{\rm cost}I_0$ 的仿射直线，截距的主项与 $N$ 无关，余项随 $N$ 增大而减小。时长边界同理。

% TODO: 下列 kappa*I0 与相对偏差来自独立核查脚本，接入 reproducibility/population.py 后以正式输出替换。
西安参数下，成本边界以及 $T_{\max}=45,150$ d 时长边界的一阶截距 $\kappa I_0$ 分别约为 $9.6\times10^{-6}$、$9.0\times10^{-6}$ 和 $2.7\times10^{-6}$ 人。在 $N\in[N_{\rm floor},\,1.316\times10^7]$ 的采样点上，三条精确边界相对于直线 $\eta=N\theta(0)$ 的最大相对偏差不超过 $5.5\times10^{-7}$。本次固定 $\theta=0.002$ 的六个人口规模测试中，$\Delta t$ 和 $J$ 的相对极差分别为 $4.71\times10^{-8}$ 和 $1.17\times10^{-7}$。这些是有限参数测试的数值结果，不构成连续人口区间上的误差认证。


% ---------------------------------------------------------------------
% 【改动 2】§9.5 第一段：原第 1616 行中的一句
%   原文：由式~\eqref{eq:s1:Imax-no}，$i_{\max}^{no}\approx0.1047$。
% ---------------------------------------------------------------------
由式~\eqref{eq:s1:Imax-no}，按全市人口计算得 $i_{\max}^{no}\approx0.1047$。


% ---------------------------------------------------------------------
% 【改动 3】§9.5 四条边界与占优区域：替换原第 1644–1671 行
%   原文开头：“峰值、成本、时长和触发条件在 $(N,\eta)$ 平面上分别为”
%   至：“……$i_{\max}^{no}N^\ast_\infty\approx9.6\times10^{3}$）。”
%   align 仍为 4 行编号，编号不变。
% ---------------------------------------------------------------------
峰值、成本、时长和触发条件在 $(N,\eta)$ 平面上的边界分别为
\begin{align}
  \text{峰值线：}\ &
  \eta=I_{\rm peak}^{\rm T}
  &&(\text{占优侧 }\eta\le I_{\rm peak}^{\rm T}),\\
  \text{成本线：}\ &
  J(\eta,N)=J^{\rm T}
  &&(\text{占优侧 }J(\eta,N)\le J^{\rm T}),\\
  \text{时长线：}\ &
  \Delta t(\eta,N)=T_{\max}
  &&(\text{占优侧 }\Delta t(\eta,N)\le T_{\max}),\\
  \text{触发线：}\ &
  \eta=I_{\max}^{no}(N)
  &&(\text{可行侧 }\eta<I_{\max}^{no}(N)),
\end{align}
其中 $J(\eta,N)$、$\Delta t(\eta,N)$ 和 $I_{\max}^{no}(N)$ 均按固定绝对 $I_0$ 与 $S_0=N-I_0$ 计算。占优区域按这些实际指标定义为
\[
\mathcal{W}_{\rm pcd}
=
\bigl\{
(N,\eta):
I_0<\eta<I_{\max}^{no}(N),\
\eta\le I_{\rm peak}^{\rm T},\
J(\eta,N)\le J^{\rm T},\
\Delta t(\eta,N)\le T_{\max}
\bigr\},
\]
不另设时长上限时去掉最后一项。峰值线是精确的水平线；由第\ref{sec:dom:setup}节的一阶近似，成本线和时长线分别接近直线 $\eta=\theta_{\rm cost}N$ 与 $\eta=\theta_{\rm dur}(T_{\max})N$，触发线接近 $\eta=i_{\max}^{no}N$。图~\ref{fig:dom} 按这些直线绘制，阈值比例取全市人口 $N=1.316\times10^7$ 下的计算值。直线只用于显示集合形状；涉及集合归属或交点的结论均直接由 $J$ 和 $\Delta t$ 判断。

比较区域的下界由成本与时长条件共同确定，其人口上限对应 $\eta=I_{\rm peak}^{\rm T}$ 上的交点 $(N^\ast(T_{\max}),I_{\rm peak}^{\rm T})$。该点仍满足严格触发条件（$152.55$ 远小于 $I_{\max}^{no}(N^\ast_\infty)\approx9.6\times10^{3}$）。


% ---------------------------------------------------------------------
% 【改动 4】图 fig:dom 的图注：原第 1687 行中的一处
%   原文：(a) 峰值、成本、时长和触发条件；人口下界右侧……
% ---------------------------------------------------------------------
(a) 峰值、成本、时长和触发条件，后三者按全市人口下的阈值比例画为直线；人口下界右侧……


% ---------------------------------------------------------------------
% 【改动 5】主人口区间：替换原第 1692–1697 行
%   原文开头：“仅考虑峰值和成本时，条件 $\theta_{\rm cost}N\le\eta\le I_{\rm peak}^{\rm T}$ 与人口下界的交集为”
%   至：“\end{equation}”。其后“在此范围内，存在阈值使……”一段保留不动。
%   区间之间的等号改为 \approx，标签 eq:dom:main-interval 不变。
% ---------------------------------------------------------------------
仅考虑峰值和成本时，由于固定 $N$ 时 $J$ 随 $\eta$ 严格递减，存在满足 $\eta\le I_{\rm peak}^{\rm T}$ 与 $J(\eta,N)\le J^{\rm T}$ 的阈值，当且仅当 $J(I_{\rm peak}^{\rm T},N)\le J^{\rm T}$。固定绝对初值时，人口上界 $N^\ast_\infty$ 取为
\[
J\bigl(I_{\rm peak}^{\rm T},N^\ast_\infty\bigr)=J^{\rm T}
\]
在本节计算范围内的根；固定归一化初值时，它与式~\eqref{eq:dom:Nstar-infty} 一致。表~\ref{tab:dom:thresholds} 中的 $N^\ast(T_{\max})$ 同样由 $\eta=I_{\rm peak}^{\rm T}$ 上的 $J$ 和 $\Delta t$ 直接求得。与人口下界取交后，有
\begin{equation}
  N_{\rm eff}\in\bigl[\,N_{\rm floor},\,N^\ast_\infty\,\bigr]
  \approx\bigl[\,1.06\times10^{4},\ 9.22\times10^{4}\,\bigr],
  \label{eq:dom:main-interval}
\end{equation}
